
\subsection{全ての滑らかな四次del PezzoにおけるJacobian枝の相対剛性}
\label{subsec:dp-universal-jacobian-branch}

定理\ref{thm:dp-rigidity}は一般の$(S,[\sigma])$の
\emph{全延長スキーム}の一意性を与える。
次の定理は全ての滑らかな$S$と非零切断に対し、
標準Jacobian型延長の\emph{局所}剛性を与える。
特に、離れた階数四解の不存在は結論しない。

\begin{lemma}[Jacobian構造と可換な定数双ベクトルの消滅]
\label{lem:dp-no-compatible-constant-bivector}
$S=V(q_0,q_1)\subset\PP^4$を任意の滑らかな四次del Pezzo、
$0\ne h\in H^0(\PP^4,\cO(1))$とし、
\[
 J_h=\tfrac14J(q_0,q_1,h),\qquad
 \iota_{J_h}\Omega=\tfrac14dq_0\wedge dq_1\wedge dh
\]
と置く。このとき、定数双ベクトル$A\in\wedge^2\CC^5$について
\[
 \boxed{\quad [J_h,A]=0\quad\Longrightarrow\quad A=0.\quad}
\]
Schouten括弧は斉次アフィン錐上のstrictな括弧である。
\end{lemma}
\begin{proof}
滑らかな二次式pencilを同時対角化し、
\[
 q_0=\sum_{i=0}^4x_i^2,\quad
 q_1=\sum_{i=0}^4\lambda_ix_i^2,\quad
 h=\sum_{i=0}^4c_ix_i,\qquad
 \lambda_i\ne\lambda_j\ (i\ne j)
\]
と書く。$c=(c_i)\ne0$である。
$A=(A_{ij})$を反対称行列とし、
$d_i=(Ac)_i=\sum_jA_{ij}c_j$、
$M=A\Lambda-\Lambda A$、
$\Lambda=\operatorname{diag}(\lambda_i)$とする。
$M$は対称行列で
$M_{ij}=(\lambda_j-\lambda_i)A_{ij}$である。
\[
 F=A(dq_0,dq_1)=4\sum_{i<j}
     (\lambda_j-\lambda_i)A_{ij}x_ix_j,\quad
 U=A(dq_0,dh)=2\sum_i d_ix_i,\quad
 V=A(dq_1,dh)=2\sum_i\lambda_id_ix_i.
\]
定数双ベクトルによる収縮を
\[
 \beta:=\iota_A(dq_0\wedge dq_1\wedge dh)
       =F\,dh-U\,dq_1+V\,dq_0
\]
と書く。$J_h$と$A$は発散零なので、
Schouten括弧と発散の標準恒等式により
$[J_h,A]=0$なら$d\beta=0$となる。

各$i<j$について$d\beta$の
$x_i\,dx_i\wedge dx_j$と$x_j\,dx_i\wedge dx_j$
の係数は、それぞれ
\begin{equation}\label{eq:dp-general-compatible-two}
 \begin{aligned}
 -4(\lambda_j-\lambda_i)(A_{ij}c_i+d_j),\\
 \phantom{-}4(\lambda_j-\lambda_i)(A_{ij}c_j-d_i).
 \end{aligned}
\end{equation}
従って全ての$i\ne j$について
\[
 d_j=A_{ji}c_i
\]
となる。$j$を固定して$i\ne j$の四式を加えると
\[
 d_j=\sum_{i\ne j}A_{ji}c_i=4d_j
\]
であり、標数零から$d_j=0$。
このため$A_{ij}c_i=A_{ij}c_j=0$が
全ての$i\ne j$について従う。

$c_k\ne0$となる$k$を一つ取ると、
全ての$i$について$A_{ki}=0$となる。
さらに$i,j,k$が相異なるとき、
$d\beta$の$dx_k\wedge dx_i$成分にある
$x_j$の係数は、
$A_{kj}=0$を用いれば符号を除いて
\[
 4(\lambda_j-\lambda_i)A_{ij}c_k
\]
となる。$\lambda_j\ne\lambda_i$、
$c_k\ne0$だから$A_{ij}=0$である。
これで全ての$A_{ij}$が零となった。
二十個の係数恒等式\eqref{eq:dp-general-compatible-two}は、
付属の独立コードでも
$\QQ[\lambda_0,\ldots,\lambda_4,c_0,\ldots,c_4,A_{ij}]$
上で検算する。
\end{proof}

\begin{theorem}[全Poisson切断に対する標準延長の被約孤立性]
\label{thm:dp-universal-jacobian-branch-rigidity}
$S\subset\PP^4$を任意の滑らかな四次del Pezzo曲面、
$0\ne\sigma_h\in H^0(S,-K_S)$を任意のPoisson構造とする。
$J_h$をその標準Jacobian延長とすると、規格化した全延長スキーム
$\ExtP_S(\sigma_h)$の$J_h$における局所環は
\[
 \boxed{\quad
 T_{J_h}\ExtP_S(\sigma_h)=0,\qquad
 \cO_{\ExtP_S(\sigma_h),J_h}\simeq\CC.\quad}
\]
従って$J_h$は被約な孤立点で、
開かつ閉な連結成分をなす。
$S$と$\sigma_h$を固定して$J_h$から出発する
非自明な形式的Poisson延長変形は存在しない。
\end{theorem}
\begin{proof}
命題\ref{prop:full20}により全双ベクトル延長は
$P=J_h+q_0A+q_1B$、
$A,B\in\wedge^2\CC^5$の二十変数で一意に表される。
$J_h$は$q_0,q_1$をstrict Casimirとするため、
$\mathcal F(P)=E\wedge\Jac(P)$の$A=B=0$における
線形化は、零でない全体尺度を除き
\[
 q_0\,E\wedge\ell_A(J_h)+q_1\,E\wedge\ell_B(J_h)
\]
である。$q_0,q_1$は互いに素で、
$U_A:=E\wedge\ell_A(J_h)$と$U_B$は
斉次二次の四重ベクトルである。
$q_0U_A+q_1U_B=0$なら
$U_A=q_1C,\ U_B=-q_0C$となる
定数四重ベクトル$C$がある。
$E\wedge U_A=0$から$E\wedge C=0$を得て、
$C=0$となる。従って$U_A=U_B=0$。

$\ell_A(J_h)$は発散零の線形係数三重ベクトルである。
五座標のEuler場$E$に対するKoszul完全性から
$E\wedge\ell_A(J_h)=0$なら
$\ell_A(J_h)=E\wedge T$となる定数双ベクトル$T$がある。
ところが$\delta(E\wedge T)=3T$であり、
$\delta\ell_A(J_h)=0$なので$T=0$。
同様に$\ell_B(J_h)=0$となる。
補題\ref{lem:dp-no-compatible-constant-bivector}から
$A=B=0$が従い、Jacobi線形化は二十変数上で単射である。

従って$J_h$におけるZariski接空間は零。
この閉点のNoether局所環の極大イデアルを
$\mathfrak m$とすると
$\mathfrak m/\mathfrak m^2=0$である。
Nakayamaの補題により$\mathfrak m=0$、
よって局所環は$\CC$となり、
この点は開かつ閉な孤立成分である。
\end{proof}

\begin{theorem}[普遍incidenceにおける標準Jacobian枝の開閉分離]
\label{thm:dp-global-standard-branch-splitting}
$B$を、滑らかな四次del Pezzo曲面$S\subset\PP^4$と
非零切断の射影類$[\sigma]\in\PP H^0(S,-K_S)$
の組をパラメータ化する既約な三十次元の空間とする。
$\mathscr I\to B$を、$\sigma$を誘導する非零ambient
Poisson構造の射影類の普遍incidenceとする。
標準Jacobian延長は切断
$s:B\to\mathscr I$を定める。
このとき
\[
 \boxed{\quad
 s(B)\subset\mathscr I\ \text{は開かつ閉な連結成分であり、}
 \qquad
 \mathscr I\simeq B\ \sqcup\ \mathscr I_{\mathrm{exc}}
 \quad}
\]
という$B$上のスキームの分解がある。

特に、滑らかな$S$と非零の$\sigma$を動かしながら
例外的なPoisson延長が標準Jacobian延長へ
合流することはない。
ここで$\mathscr I_{\mathrm{exc}}$は空でもよく、
その空性と階数四部分の存在は未解決である。
\end{theorem}
\begin{proof}
$B$は滑らかな二次式pencilを表す
$\operatorname{Gr}(2,\operatorname{Sym}^2V^*)$の
26次元開集合の上に、五次元の反標準切断空間の
射影$\PP^4$を取るので既約かつ三十次元である。
$\mathscr I\to B$は、非零制限による尺度規格化のもとで
局所的に二十次元アフィン空間束内のJacobi方程式で
定義される有限表示の分離射である。
$J_h$は$pencil$の基底変更で射影尺度しか変わらないため、
標準延長は大域切断$s$を与える。

前定理から全ての幾何点$b\in B$で
$s(b)$の相対接空間は零である。
相対微分の有限生成性とNakayamaの補題により、
$\mathscr I\to B$は$s(B)$の近傍でunramifiedとなる。
unramified射の切断は開埋込み、
分離射の切断は閉埋込みなので、
$s$は$\mathscr I$の開かつ閉な埋込みである
\cite[Proposition 41.6.1]{StacksUnramifiedSection}。
$B$は既約ゆえ連結であり、
その像は一つの連結成分となる。
残る補集合を$\mathscr I_{\mathrm{exc}}$と置けば主張を得る。
\end{proof}

\begin{corollary}[判別式の五つの固有超平面切断]
\label{cor:dp-eigen-hyperplane-rigidity}
各滑らかな四次del Pezzoの二次式pencilは、
五つの単純判別式零点に対応する五つの
固有超平面$H_i$を持つ。
それらによる反標準曲線$S\cap H_i$は
いずれも滑らかな楕円四次曲線であり、
標準Jacobian延長は上定理の被約孤立点である。

この五つの特別な切断については、
補題\ref{lem:dp-no-compatible-constant-bivector}
の定数双ベクトル線形化の十次小行列式も明示的に
求められる。$h=x_0$、$\lambda_0=0$の同時対角座標で、
$(01,02,03,04,12,13,14,23,24,34)$を列とし、
行を
\[
(012;3),(012;4),(013;4),(023;4),
(123;1),(123;2),(123;3),(124;1),(124;2),(134;1)
\]
と取ると、行列式は
\[
 -(\lambda_1-\lambda_2)(\lambda_1-\lambda_3)
 (\lambda_1-\lambda_4)^2(\lambda_2-\lambda_3)
 (\lambda_2-\lambda_4)^2(\lambda_3-\lambda_4)^3\ne0.
\]
これは独立の
\texttt{check\_dp\_eigen\_rigidity.py}でも
$\QQ(\lambda_1,\ldots,\lambda_4)$上で検算されている。
\end{corollary}

\begin{remark}[普遍剛性と混同しない]
全ての標準Jacobian延長は相対incidenceの一つの
開閉成分を形成するが、これだけで
$\mathscr I_{\mathrm{exc}}=\varnothing$は従わない。
特殊な$(S,[\sigma])$の例外的階数四延長が存在するなら、
標準枝とは切り離された別の連結成分に生じる。
従って今後は標準枝の無限小特異性を探索するのではなく、
独立の例外的解をPoisson零スキーム、余法線Lie代数と
退化五次因子から検出する必要がある。
\end{remark}
